\documentclass[../main.tex]{subfiles}

\begin{document}
\begin{problema}
	Demuestre que la ecuación

	\begin{equation*}
		A \abs{z}^{2} + Bz + \overline{B}\overline{z} + C = 0,
	\end{equation*}

	con \(A,C\in \mathbb{R}\) y \(B\in \mathbb{C}\), no todos cero, describe una recta o una
	circunferencia en \(\mathbb{C}\).

	\begin{proof}
		Sean \(A, C \in \mathbb{R}\) y \(B = B_{1} + iB_{2}\in \mathbb{C}\), entonces

		\begin{equation*}
			A \abs{z}^{2} + Bz + \overline{Bz} + C = 0
		\end{equation*}

		pues \(\overline{Bz} = \overline{B}\overline{z}\) y además observamos que

		\begin{equation*}
			\Re(z) = \dfrac{z + \overline{z}}{2}. \quad \therefore\; 2 \Re(Bz) =  Bz + \overline{Bz}.
		\end{equation*}

		Así,

		\begin{equation}
			A \abs{z}^{2} + 2\Re(Bz) + C = 0,
			\label{eq:gral-eq}
		\end{equation}

		\begin{itemize}
			\item Si \(A = 0\),

			      \begin{equation*}
				      Bz + \overline{Bz} + C = 0.
			      \end{equation*}

			      Por lo tanto, la ecuación describe una recta.
			\item Si \(A \neq 0\),

			      Desarrollando la \zcref{eq:gral-eq},

			      \begin{align*}
				      A \bigl(x^{2} + y^{2}\bigr) + 2B_{1}x - 2B_{2}y + C                         & = 0,              \\
				      x^{2} + y^{2} + 2\dfrac{B_{1}}{A}x - 2\dfrac{B_{2}}{A}y                     & = - \dfrac{C}{A}, \\
				      \Bigl(x + \tfrac{B_{1}}{A}\Bigr)^{2} + \Bigl(y - \tfrac{B_{2}}{A}\Bigr)^{2} & =
				      \dfrac{B_{1}^{2} + B_{2}^{2} - CA}{A^{2}},                                                      \\
				      \Bigl(x + \tfrac{B_{1}}{A}\Bigr)^{2} + \Bigl(y - \tfrac{B_{2}}{A}\Bigr)^{2} & =
				      \dfrac{\abs{B}^{2} - CA}{A^{2}}.
			      \end{align*}

			      Por lo tanto, la ecuación describe una circunferencia si \(\abs{B}^{2} - CA \neq 0\).
		\end{itemize}
	\end{proof}
\end{problema}
\end{document}
